% GENERATED_BY: oc_core_monograph_translation_draft_pipeline_v1.py
% AI_ASSISTED_TRANSLATION_DRAFT: TRUE
% THEOREM_REVIEW_STATUS: PENDING
% THEOREM_REVIEW_MODE: FORMAL_AI_GATE__CODEX_CLI_CHATGPT
% THEOREM_REVIEW_REVIEWER_ID:
% THEOREM_REVIEW_AUTHORITY_CLASS:
% THEOREM_REVIEWED_AT:
% THEOREM_REVIEW_DOSSIER_REF:
% SOURCE_LANGUAGE: EN
% TARGET_LANGUAGE: RU
% SOURCE_EN_REF: appendix/D_oc_core_1_3_source_audit_appendix.tex
\section{Core 1.3.3 Provenance and Corpus Appendix}

This appendix records the principles by which the 1.3 release remains source-first while still
operating as a публичныйation program rather than as a static archive.

\subsection{Why an appendix is still needed}

The main body of the monograph is written to be read.
The source-audit machinery is written to be checked.
Those are related but non-identical tasks.
If the full provenance burden were pushed into the main body, the book would become harder to read.
If it were omitted entirely, the reader would need to trust a hidden workflow.
The appendix therefore carries the discipline without replacing the main narrative.

\subsection{Witness classes}

The release uses several witness classes:

\begin{itemize}
\item the публичный LaTeX corpus as the исходно-авторские textual witness,
\item the archival compiled PDF as the stable page-and-line witness,
\item the bounded journal-core extraction as the outward review-focused derivative witness,
\item and companion packets as support witnesses for chronology, теорему fate, and release
      packaging.
\end{itemize}

None of these witnesses should be mistaken for the others.
The свидетельство источника is not the same as the release manifest.
The release manifest is not the same as the теорему.
The journal-core статья is not the same as the monograph.
The 1.3 публичныйation line remains healthy only when these roles stay distinct.

\subsection{Questions that the appendix answers}

This appendix is meant to answer concrete review questions.

\begin{itemize}
\item From which публичный source stack is the present volume built?
\item Which outward статья is derived from it?
\item Why is the monograph larger than the journal-core статья?
\item What justifies calling the volume 1.3 rather than simply republishing 1.2?
\item How should later the methodological staging framework and channel projections interpret the distinction between formal
      master, review статья, and пакет теоремной опоры?
\end{itemize}

The release is stronger when those questions can be answered inside the volume rather than by
pointing the reader to an unrelated dashboard.

\subsection{Current boundary}

Core 1.3.3 should therefore be read as a source-grounded clarification and clarification of the formal
публичныйation shell.
The monograph does not pretend that every future projection is already proved.
The journal core does not pretend to be the whole theory.
The companion packets do not pretend to be the manuscript.
That three-way distinction is the governing boundary of the current release.

\clearpage
